\documentclass[../main]{subfiles}
\begin{document}
\ifSubfilesClassLoaded{\setcounter{page}{7}}{}
\section{Передмова}
\label{sec:02_predislovie}

Критикуючи ґотську програму, Маркс писав: «Ми маємо справу не з таким комуністичним суспільством, яке розвинулося на власній основі, а, навпаки, з таким, яке щойно виходить із капіталістичного суспільства і яке тому, в усіх відношеннях, в економічному, моральному та інтелектуальному плані, зберігає ще залишки старого суспільства, з надр якого воно виросло». І далі: «На вищій фазі комуністичного суспільства, після того як зникне поневолююче людину підпорядкування її поділу праці; коли зникне разом з цим протилежність між інтелектуальною та фізичною працею; коли праця перестане бути лише засобом для життя, а стане самою першою потребою життя; коли разом із всебічним розвитком індивідів зростуть і продуктивні сили, і всі джерела суспільного багатства почнуть повноцінно функціонувати, лише тоді можна буде остаточно подолати вузький горизонт буржуазного права, і суспільство зможе написати на своєму прапорі: Кожному – за здібностями, кожному – за потребами!

Ідеї приватної власності та панування капіталу протиставляється ідея комунізму. Реальності, де панує капітал, можна протиставити лише комуністичне суспільство. Але що таке комунізм?

У цій брошурі нас цікавить те, що Маркс назвав найвищою фазою комунізму, тобто розвинене на власній основі комуністичне суспільство, або комунізм як такий. Нижча фаза комунізму, або перехідний період від капіталізму до комунізму, який пізніше історично закріпився під назвою «соціалізм», буде розглянута нами пізніше. Зараз йдеться про те, до чого саме має відбутися цей перехід, або про те, чого хочуть досягти комуністи в результаті революційних перетворень суспільного життя.

Нас тут цікавлять ті визначальні характеристики, які роблять суспільство саме комуністичним і відрізняють його від усіх попередніх комунізму форм виробництва та організації суспільного життя. Тому ми поки що відкладемо опис технологій при комунізмі. Вони можуть бути різними (і, безумовно, швидко змінюватимуться). Наприклад, капіталізм за часів панування парового двигуна та капіталізм із сучасною цифровою технікою дуже сильно відрізняються один від одного. Але те, що робить капіталізм капіталізмом, суспільством розвиненого товарного виробництва, де товаром є, зокрема, і робоча сила, а капітал – домінуючою суспільною формою відносин, є характерним, необхідним і незмінним для капіталізму взагалі, як у XIX, так і в XXI столітті. Розуміння цих необхідних моментів у їхньому необхідному взаємозв’язку і становить поняття капіталізму.

Так само й у поняття комунізму, як і будь-якого іншого явища, входить лише те, що є необхідним і достатнім у взаємопов’язаних аспектах, щоб суспільство було комуністичним.

Складність полягає в тому, що ми маємо справу з явищем, яке ще не стало реальністю. Історія ще не знала (принаймні, на нашій планеті) комуністичного суспільства, яке розвивалося б на власній основі (первісний комунізм у цьому випадку не береться до уваги).

Проте, комунізм вже давно став наукою. Увесь хід історичного розвитку думки породив ідею комунізму, а хід суспільного життя, принаймні, постійно демонструє межі капіталізму та необхідність переходу до нової основи суспільного розвитку. Капіталістичні кризи надвиробництва супроводжуються колосальним знищенням не лише вироблених товарів чи окремих підприємств, але й часто цілих колись процвітаючих мегаполісів.

Якщо розглядати суспільне виробництво як спосіб взаємодії людини з природою, стає очевидно, що воно є нераціональним і завдає нашій планеті шкоди таких масштабів, що це може стати непоправною проблемою. Виробництво і надмірне виробництво тієї ж самої упаковки, зайвої та безглуздої з точки зору розумної взаємодії з природою, але необхідної для того, щоб товар був проданий, – завдає величезної шкоди. Виробництво самих товарів, які не збагачують і не розвивають, а часто й калічать людину, але поглинають велику кількість ресурсів, забруднюють ґрунт, воду та повітря – ось реальність капіталізму. Відсутність планомірної в масштабах усього суспільства, продуманої переробки відходів, а також розумного розподілу та використання всіх природних ресурсів і виробничих можливостей. Все це створює загрозу незворотної екологічної катастрофи.

Паралельно з таким рукотворним руйнуванням природи, що нас оточує, відбувається деградація людини. Це не лише абсолютна бідність, в яку занурена значна частина працюючого населення Землі: люди голодують, живуть в антисанітарних умовах, не мають доступу до освіти та медицини; але йдеться також про відносне зубожіння, коли створюються нові потреби, задоволення яких стає необхідністю навіть для того, щоб існувати як продавець власної робочої сили, але засобів для їх задоволення недостатньо. Цю тенденцію ще зазначав Енгельс, але в наші дні, коли капіталізм не може існувати без цілеспрямованого виробництва потреб, вона багатократно посилюється і набуває загрозливих масштабів, особливо на периферії та напівпериферії капіталістичного світу, де в суспільному масштабі робоча сила не відновлюється в повній мірі. Це підриває самі основи капіталістичного виробництва. Але це ще не все.

Щоб відтворюватися в розширеному масштабі, сучасний капітал відроджує та підпорядковує собі докапіталістичні способи експлуатації людини людиною. Рабство та кріпацтво в сучасному світі існують доти, поки вони виявляються вигідними та сприяють самозростанню вартості. Наприклад, за даними МОП, опублікованими в 2014 році, прибутки від використання рабської праці в усьому світі перевищують 150,2 мільярди доларів (як бачите, вигода від рабства обчислюється і може обчислюватися суто в капіталістичний спосіб, а не в натуральній формі, як це було в епоху рабовласництва).

Капітал постійно потребує простору для самозростання, а для цього йому потрібно витіснити інший капітал з уже наявних ринків, перерозподілити ресурси на свою користь. Це породжує конфлікти не лише всередині держави, але, головним чином, між ними, що визначає локальні військові дії та веде до чергової світової війни. А коли ми маємо тенденцію (підготовку, розгортання, окремі спалахи) світової війни та ядерну зброю, які тільки непередбачувані випадковості можливі! При цьому сама війна дає потужний поштовх до розвитку широкого, масового, стихійного антикапіталістичного руху. Наскільки послідовною буде ця лінія, і чи вона вестиме до реального виходу за межі капіталізму – сказати поки що важко. Взагалі, гарантії того, що людство переживе капіталізм, ніхто зараз дати не може. Єдине, про що це свідчить, – це про межу розвитку суспільства, де панує капітал.

Те, що капіталізм колись припинить своє існування, незаперечно. Питання лише в тому, чи закінчиться він разом із людством, чи людство навчиться жити інакше. І як саме інакше можна жити, щоб і людство в цілому, і всі його різноманітні складові розвивалися, не заважаючи, а допомагаючи одне одному? Наукова теорія комунізму дає відповідь на це, зовсім не застаріле, а надзвичайно актуальне для людства питання.

Основні моменти цієї теорії ми й спробуємо тут викласти. Її створили Маркс і Енгельс. У розвиток розуміння комунізму свій внесок зробили й пізніші мислителі-марксисти. Серед них особливої уваги заслуговує теоретична праця Леніна.

Значення термінів «комунізм» і «соціалізм» не сформувалося одразу. Тривалий час вони вживалися як синоніми, або комунізм навіть могли розуміти як попередній етап соціалізму. Тому слід однаковою мірою звертатися як до праць Маркса та Енгельса, де вживається термін «комунізм», так і до тих, де йдеться про «соціалізм», якщо за змістом викладеного зрозуміло, що йдеться про комуністичне суспільство, про комуністичний рух або про теорію комунізму. Що стосується Леніна та інших марксистів ХХ століття, то в їхніх працях, як і в пізніх працях Маркса, термінологія вже більш-менш усталилася. У статті «Великий початок» Ленін спеціально нагадує читачам, що «наукова відмінність між соціалізмом і комунізмом полягає лише в тому, що перше слово означає першу ступінь нового суспільства, яке виростає з капіталізму, а друге слово — вищу, наступну ступінь». Там, де ми використовуватимемо ці слова без посилань і особливих пояснень, їх слід розуміти відповідно до сучасного вживання.

Щодо аргументу проти комунізму, який з різних боків повторює всяка і вся з лав ідеологічних прислужників сучасного правлячого класу капіталістів: «комунізм – це утопія, ми вже намагалися побудувати комунізм, і нічого (або нічого доброго) з цього не вийшло», – на це ми відповімо наступне. Так, комуністичне суспільство побудувати поки що не вдалося, але стверджувати, що нічого доброго з цього не вийшло, – це просто брехня. Успіхи соціалістичного Радянського Союзу свого часу не могли заперечувати навіть найзатятіші антикомуністи, навіть у США. Ті, хто пам’ятає недавню історію ХХ століття, добре знають, що рух суспільства до комунізму дозволив народам, які жили в СРСР, досягти небувалих раніше, за часів феодалізму та капіталізму, висот, і при цьому за найкоротший термін. Про це свідчать як сухі цифри та статистичні дані, так і досягнення в галузі культури, які досі не мають собі рівних.

Якщо порівняти перші десятиліття існування суспільства, коли спостерігався рух у напрямку до комунізму, і той самий період розпаду СРСР, коли країни «незалежних держав» почали розходитися різними шляхами, то ми побачимо пряму залежність між рухом до комунізму та розвитком усього суспільства. Йдеться не лише і не стільки про промисловість і виробництво матеріальних благ, скільки про можливості для розвитку людини та людських здібностей, які були доступні широким верствам населення.

Вже у 1925 році в початкових школах навчалося на 2 250 000 учнів більше, ніж у 1913 році. Кількість учнів у професійно-технічних школах зросла вдвічі. З громадянської війни суспільство одразу перейшло до навчання. Право на освіту було здобуто для тих, хто до громадянської війни був приречений на невігластво. Задумайтеся: Радянському Союзу лише рік, а у нього вже є чим похвалитися перед капіталістичним світом, де в освіті домінує класовий принцип. Якісна освіта в капіталістичному світі була і залишається недоступною для більшості.

За 20 років після Жовтневої революції вдалося вирішити проблему безпритульності, і дітям, які втратили батьків під час Першої світової та Громадянської війни, давали не просто дах над головою та їжу, а й освіту, професію, можливість почати життя. Водночас, на уламках колишньої єдиної соціалістичної країни знову з’явилися безпритульні діти, і в більшості випадків їхні батьки живі. Рівень загальної освіченості населення також знизився і продовжує знижуватися. Системи освіти колишніх радянських республік прагнуть відповідати новим соціальним реаліям: все більш вузькій спеціалізації та поглибленню розшарування в освітніх можливостях для дітей з бідних і заможних сімей. Безліч красномовних свідчень про цей процес ви знайдете у фільмі Костянтина Сьоміна та Євгена Спіцина «Останній дзвінок».

У СРСР програма ліквідації неписьменності (лікбез – з 1920 року) та загального початкового навчання (всеобуч – з 1930 року) змогла вирішити проблеми початкової освіти менш ніж за двадцять років – і це в аграрній країні, де до 1917 року було понад 80% неписьменного та малописьменного населення. До 1936 року в СРСР було навчено близько 40 мільйонів неписьменних. У 1933–1937 роках лише в облікованих школах лікбезу займалися понад 20 мільйонів неписьменних і близько 20 мільйонів малописьменних. За даними перепису 1939 року, грамотність осіб віком від 16 до 50 років наближалася до 90%. Саме ці люди, а не вигадані герої з хрестами на грудях і Богом у серці (на чому наполягав фільм Микити Михалкова «Цитадель» та інші зразки сучасної пропагандистської індустрії), перемогли найпотужнішу, найсучаснішу та найозброєнішу на той момент військову машину нацистської Німеччини, на яку до моменту вторгнення в СРСР працювала більша частина європейської промисловості.

Отже, за 20 років було практично ліквідовано неписьменність, створено мережу медичних установ, де кожен мав доступ до медичної допомоги (лікарі та вчителі з’явилися там, де їх раніше ніколи не було), з’явилися школи, бібліотеки, заклади культури, було надано потужний поштовх для розвитку науки, зокрема й завдяки появі відповідної інфраструктури. Пріоритетне всебічне розбудовування комунізму, спрямоване на розвиток людини (освіта, охорона здоров’я, наука, культура), дозволило перемогти фашизм, який загрожував усьому світу, а після важкої руйнівної війни – у найкоротші терміни відбудувати країну.

Що ж принесло життя колишнім радянським республікам після розпаду СРСР – повне згортання комуністичних тенденцій у розвитку суспільства?

Лише в Україні після розпаду СРСР, за даними Державного комітету статистики до 2014 року, абсолютна втрата населення становила майже 6,5 мільйона осіб (до розпаду спостерігалося позитивне зростання). Ця цифра не враховує тих, хто «числиться» громадянами України, але був змушений покинути свій дім і поїхати на заробітки. Це саме по собі свідчить про колосальну соціальну катастрофу. Звісно, тут згадають голод 1932-1933 років. Масовий голод – завжди жахлива трагедія, незалежно від її справжніх причин. Але навіть за найантибільшовицькими підрахунками (саме підрахунками, а не безпідставними ідеологічними заявами), проведеними Інститутом демографії НАН України, голод в Україні забрав не 7 і не 10, а 3,9 мільйона людських життів. Враховувалися не лише втрати, а й порівнювалася чисельність населення з прогнозованою чисельністю з урахуванням значного позитивного приросту в роки до і після голоду. Якщо за аналогічною методикою обчислити абсолютні втрати населення України після розпаду СРСР, то картина буде набагато жахливішою, ніж та, що отрималася в результаті віднімання майже 6,5 мільйона.

У Литві, Латвії та Естонії також спостерігається значне скорочення населення. У РФ за рахунок міграції зафіксовано меншу втрату населення: зі 148,6 мільйона у 1991 році до 144,3 мільйона у 2016 році.

Якщо врахувати трудову міграцію в інших напрямках і структуру зайнятості в самих середньоазіатських республіках, стане зрозуміло, що загальна якість життя та можливість громадян цих республік знайти роботу, яка дозволяє утримувати сім’ю, у себе вдома значно знизилися.

Війни на території колишнього СРСР стали поширеним і звичайним явищем: конфлікт між Азербайджаном і Вірменією в Нагірному Карабаху (який розпочався ще за часів СРСР, але, по суті, розгорнувся під час розпаду), громадянська війна в Таджикистані, Придністров’я, дві чеченські війни, Грузія та Абхазія, Південна Осетія, Донбас. І в усіх цих війнах люди вбивали одне одного, зрештою, заради конкуруючих угруповань капіталу, тобто своїх власних експлуататорів.

На це можуть заперечити, що громадянська війна, яка почалася в 1917 році, також була кривавою і, за різними підрахунками, забрала від 8 до 13 мільйонів життів з обох сторін. Сюди входять не лише ті, хто загинув у бойових діях, а й ті, хто помер від хвороб, голоду та терору. Так, це правда, – відповімо ми. І, звісно, було б краще, якби війни не було. Але, по-перше, її розв’язали не більшовики і не пролетаріат, а представники старих експлуататорських класів, які не хотіли втрачати свою владу. По-друге, у цій війні трудящі відстоювали власні інтереси, а не інтереси тих, хто живе за рахунок їхньої праці. По-третє, ця війна була боротьбою за можливість здійснити кардинальні зміни в усьому суспільстві, а не просто за переділ сфер впливу капіталу. Про Велику Вітчизняну війну має сенс говорити окремо, оскільки це була боротьба з фашизмом, де радянський народ відстоював не лише власні інтереси, але й інтереси всього людства.

Як бачите, навіть таке просте порівняння перших десятиліть існування СРСР і життя республік після його розпаду дає можливість говорити про переваги суспільства, де є простір для тенденції руху до комунізму. І це при тому, що процес відбувався зовсім не гладко, і було багато того, що аж ніяк не можна назвати комуністичним. Йдеться про чисто капіталістичні процеси та явища в соціалістичному суспільстві. Саме про них так багато кричать критики СРСР, незважаючи на те, що ця критика, там, де вона цілком справедлива, виявляється самокритикою щодо капіталістичного світу.

Тенденція до руху в бік комунізму в СРСР дозволила всього через дванадцять років після Великої Вітчизняної війни, яка забрала мільйони життів і зруйнувала значну частину країни, запустити перший штучний супутник Землі. А незабаром і відправити першу людину в космос. Будівництво комунізму дозволило розробити в науковому плані ті ідеї, які для нас, у нашому нинішньому стані суспільства, скоріше є ідеями майбутнього, ніж минулого, оскільки їх реалізація несумісна з приватною власністю на засоби виробництва. Цей список можна продовжувати довго і включати в нього не тільки радянські досягнення, а й досягнення інших соціалістичних країн, – які виникли саме на тому шляху, який бачив Ленін у світлі гіпотези про перемогу світової революції.

Звісно, в цьому грандіозному проєкті були не лише досягнення, а й відступи та провали.

Але, навіть якщо не брати до уваги все вищесказане, аргумент «не вдалося побудувати комунізм, отже, це взагалі неможливо і нічого з цього не вийде» має той самий логічний недолік, що й заяви противників ідей, які, наприклад, лягли в основу виникнення та розвитку авіації. Противники стверджували, що неможливо створити літальні апарати, важчі за повітря. Це те саме, що заявити: людина чи людство не зможе чомусь навчитися лише тому, що у неї досі це не вийшло, хоча спроби й були.

Такі нелогічні аргументи не передбачають вивчення суті справи та розгляд аргументів по суті. Тому вони й не стосуються науки про комунізм. Невдача першої спроби людства досягти комунізму свідчить лише про необхідність детального вивчення історичного досвіду, щоб зрозуміти причину цієї невдачі, а також більш ретельного вивчення умов, які необхідно досягти для переходу до комунізму. Саме існування та структура комуністичного ідеалу, що виник як осмислення тенденцій і суперечностей капіталістичного суспільства, вимагають цього.

Отже, оскільки комунізм, окрім усього іншого, став наукою, давайте приступимо до його вивчення. Ця брошура, звісно, не може претендувати на детальне та розгорнуте дослідження комунізму. Скоріше, це текст, призначений для викладення основних положень науки про комунізм, розроблених у марксистській теорії. Більш глибока та детальна розробка проблеми – справа найближчого майбутнього. Ми запрошуємо наших читачів приєднатися до цієї справи, зокрема й у формі критики нашої роботи.
\end{document}
